% document formatting
\documentclass[10pt]{article}
\usepackage[utf8]{inputenc}
\usepackage[left=1in,right=1in,top=1in,bottom=1in]{geometry}
\usepackage[T1]{fontenc}
\usepackage{xcolor}

% math symbols, etc.
\usepackage{amsmath, amsfonts, amssymb, amsthm}
\usepackage{dsfont}

% lists
\usepackage{enumerate}

% images
\usepackage{graphicx} % for images
\usepackage{tikz}
\usetikzlibrary{fit}

% code blocks
\usepackage{minted, listings} 

% verbatim greek
\usepackage{alphabeta}
\DeclareMathOperator*{\argmin}{arg\,min}
\DeclareMathOperator*{\argmax}{arg\,max}
\newcommand{\dd}{\text{d}}

\graphicspath{{./assets/images/}}

\title{02-750 Week 8 \\ \large{Automation of Scientific Research}}
 
\author{Aidan Jan}

\date{\today}

\begin{document}
\maketitle

\subsection*{Background: Regression}
The goal of regression analysis is to learn a map / function, $f \::\: \mathbb{R}^d \rightarrow \mathbb{R}$
\begin{itemize}
	\item As is the case for learning classifiers, the choice of training instances impacts generalization error, thus, we want to select them actively.
\end{itemize}
Regression models (like classifiers) can be divided into two categories:
\begin{itemize}
	\item \textbf{Parametric:} the model \textbf{summarizes} the training data using a \textbf{fixed} number of parameters (independent of the size of the training data)
	\begin{itemize}
	    \item For example, linear regression: $y = x^T \beta$
    \end{itemize}
    \item \textbf{Non-parametric:} the number of parameters is \textbf{proportional} to the amount of training data, and the predictions are often made using the actual training instances
    \begin{itemize}
	    \item For example, Gaussian Process Regression, KNN regression, Regression Trees
    \end{itemize}   
\end{itemize}

\subsection*{Background: Bias-Variance Decomposition}
A model's \textbf{generalization error} can be decomposed as follows:
\[Err(x) = Bias^2 + Variance + Irreducible~Error\]
\begin{itemize}
	\item $\text{Bias} := \mathbb{E}_D [\hat{f}(x; D)] - f(x)$
	\begin{itemize}
	    \item This is the error due to the difference between the `true' function and the one we select from our chosen hypothesis class (Recall PAC learning)
    \end{itemize}
    \item $\text{Variance} := \mathbb{E}_D \left[\left(\mathbb{E}_D[hat{f}(x; D)] - \hat{f}(x; D)\right)^2\right]$
    \begin{itemize}
	    \item This is the error due to the training data being a finite sample of $P(X, Y)$.  Different training sets will produce different parameters, creating variance.
    \end{itemize}
    \item Irreducible Error: This is the error due to any noise in the training data
\end{itemize}

\subsubsection*{The bias-variance tradeoff}
\begin{itemize}
	\item Bias and variance are generally in opposition to one another
	\begin{itemize}
	    \item Complex models tend to have low bias but high variance
	    \item Simple models (i.e., those with few parameters) tend to have high bias but low variance
    \end{itemize}
    \item Some learning algorithms intentionally introduce bias in order to decrease variance
    \begin{itemize}
	    \item For example, \textbf{Lasso} and \textbf{Ridge} regression introduce bias to reduce variance
    \end{itemize}   
\end{itemize}

\subsection*{Regression Analysis}
\begin{itemize}
	\item The field of Statistics has developed a variety of techniques for learning regression models efficiently (i.e., with low sample complexity)
	\item Examples:
	\begin{itemize}
	    \item Latin Hypercube Sampling
	    \item Randomized Block Designs
	    \item Taguchi Methods
	    \item Optimal Design of Experiments (DOE)
    \end{itemize}
\end{itemize}

\section*{Optimal Design of Experiments}
A branch of statistics devoted to pre-selecting sets of \textbf{experimental conditions} that are collectively optimal \textbf{with respect to some criterion}
\begin{itemize}
	\item DOE shares many of the same objectives as Active Learning but\dots
	\begin{itemize}
	    \item Most DOE techniques are specifically for learning \textbf{parametric} regression models
	    \item Most DOE methods select \textbf{all} the experiments \textit{a priori}, not iteratively$^*$
	    \item \textbf{Some DOE methods are more than 100 years old!}
    \end{itemize}
\end{itemize}
DOE methods often make (strong) assumptions about:
\begin{itemize}
	\item The \textbf{parametric form} of the model (e.g., linear)
	\item The nature of the noise
	\item The test density, $P(X)$
\end{itemize}

\subsection*{Example: DOE for Linear Regression}
Model (in matrix form for $k$ training instances): $\mathbf{y} = X\beta + \epsilon$
\begin{itemize}
	\item $\mathbf{y}$ is a $k \times 1$ vector containing $k \geq 1$ observations / samples
	\item $\mathbf{X}$ is the $k \times b$ \textbf{design matrix}.  Here, $b$ is the number of parameters.
	\item $\beta$ is the $b \times 1$ parameter vector
	\item $\epsilon$ is a $k \times 1$ noise / error vector
	\begin{itemize}
	    \item Assumptions
	    \begin{itemize}
	        \item $\mathbb{E}[\epsilon_i] = 0$.  (Expected value of noise is zero)
	        \item $Var(\epsilon) = \sigma^2 < \infty$  (Noise is constant and has finite variance)
	        \item $\mathbb{E}[\epsilon_i \epsilon_j] = 0$  (Noise is uncorrelated across instances)
        \end{itemize}
    \end{itemize}
\end{itemize}
The standard way of estimating $\beta$ via \textbf{ordinary least squares} (OLS)
\begin{itemize}
	\item Given a vector of $k \geq b$ observations / samples the OLS estimate is:
	\[\hat{\beta} = (X' X)^{-1} X' y\]
    \begin{itemize}
	    \item Once we estimate $\hat{\beta}$, the predicted value for any $x$ is: $\hat{y} = x\hat{\beta}$
    \end{itemize}
    \item It can be shown (via the Gauss-Markov Theorem) that:
    \begin{itemize}
	    \item OLS produces a $\hat{\beta}$ that minimizes $\sum_{i = 1}^k (\hat{y}_i - y_i)^2$
	    \item $\mathbb{E}[\hat{\beta} - \beta_{true}] = 0$  (i.e., OLS is an unbiased estimator)
	    \item $Var(\hat{\beta}) = \sigma^2 (X'X)^{-1}$
	    \item $Var(\hat{y}) = \sigma^2 x (X' X)^{-1} x'$
    \end{itemize}
\end{itemize}
Recall the bias-variance decomposition:
\[Err(x) = Bias^2 + Variance + Irreducible~Error\]
\begin{itemize}
	\item OLS is unbiased, so $Err(x) = Variance + Irreducible~Error$
	\item We know that $Var(\hat{\beta}) = \sigma^2(X'X)^{-1}$
	\begin{itemize}
	    \item $\sigma^2$ is usually unknown, but it is a constant, so we can ignore 
        \item So the \textbf{design matrix}, $X$ is really the \textbf{only thing that matters}, not the $y$'s.
    \end{itemize}
\end{itemize}
DOE methods seek to minimize the variance in the estimate (and thus generalization error) by selecting an optimal set of $k$ training instances (i.e., rows of $X$)
\begin{itemize}
	\item Notice that these instances can be selected without knowing \textbf{any} labels!
	\item Different DOE methods use different definitions for what it means to be optimal
\end{itemize}

\subsection*{Quantifying Uncertainty for Linear Regression}
\begin{itemize}
	\item Given $D_L = \{(x_1, y_1), \cdots, (x_n, y_n)\}$, the sum of squared vertical deviations from the point to the line is
	\[f(b_0, b_1) = \sum_{i = 1}^n [y_i - (b_0 + b_1 x_i)]^2\]
    \item The point estimates of $\beta_0$ and $\beta_1$, denoted by $\hat{\beta}_0$ and $\hat{\beta}_1$, are the \textbf{least squares estimates}, that is, the values that minimize $f(b_0, b_1)$
\end{itemize}
\begin{align*}
    b_1 &= \hat{\beta}_1 = \frac{\sum_{i = 1}^n (x_i - \bar{x})(y_i - \bar{y})}{\sum_{i = 1}^n (x_i - \bar{x})^2} = \frac{s_{xy}}{s_{xx}} \\
    b_0 &= \hat{\beta}_0 = \frac{\sum_{i = 1}^n y_i - \hat{\beta}_1 \sum_{i = 1}^n x_i}{n} = \bar{y} - \hat{\beta}_1 \bar{x} \\
    s_{xy} &= \sum_{i = 1}^n x_i y_i - \frac{\sum_{i = 1}^n x_i \sum_{i = 1}^n y_i}{n} \\
    s_{xx} &= \sum_{i = 1}^n x_i^2 - \frac{\left(\sum_{i = 1}^n x_i\right)^2}{n} \\
    \hat{\sigma}^2 &= s^2 = \frac{\sum_{i = 1}^n (y_i - \hat{y}_i)^2}{n - 2} = \frac{1}{n - 2} \sum_{i = 1}^n \hat{e}_i^2
\end{align*}

\subsection*{Optimality Criteria: G-Optimality}
\[\min_{x_i, i = 1:k} \max_{x \in \mathcal{X}} \text{Var}(\hat{y}_x)\]
\begin{itemize}
	\item This criterion selects the $k$ instances that minimize the maximum variance of any prediction
	\begin{itemize}
	    \item This is equivalent to minimizing the maximum element in the diagonal matrix $X(X'X)^{-1} X'$
    \end{itemize}
    \item This is the criterion proposed by Smith
\end{itemize}

\subsection*{Optimality Criteria: D-Optimality}
\[\max_{x_i, i = 1:k} |X'X|\]
\begin{itemize}
	\item Select the $k$ instances that maximize the \textit{determinant} of $X'X$
	\item Geometric Interpretation
	\begin{itemize}
	    \item $|X'X|$ is inversely proportional to the product of the axes of the \textbf{confidence ellipsoid} around the point estimate of our parameter vector $\beta$
	    \item Bigger $|X'X|$ means smaller ellipsoid, thus, higher confidence in $\beta$
    \end{itemize}
\end{itemize}

\subsection*{Optimality Criteria: A-Optimality and E-Optimality}
\textbf{A-Optimality:} (Chernoff 1953)
\[\min_{x_i, i = 1:k} \text{Tr}\left((X'X)^{-1}\right)\]
\begin{itemize}
	\item Select the $k$ instances that minimize the \textit{trace} of the inverse $X'X$
	\item Minimizes the average variance of the parameter estimates
\end{itemize}
\textbf{E-Optimality:} (Ehrenfeld 1953)
\begin{itemize}
	\item Select the $k$ instances that maximizes the minimum eigenvalue of $X'X$
	\item Minimizes the maximum variance of $\beta$
\end{itemize}

\subsection*{DOE for Nonlinear Models}
If the ``true'' model is non-linear, but can be expressed as a linear \textbf{basis function}, we can still use the previous objectives to find optimal designs
\begin{itemize}
	\item A linear basis function has the form: $f(x) = \sum_{i =1 }^b \alpha_i \phi_i(x)$
	\begin{itemize}
	    \item Here, each $\phi_i$ is one of the basis functions
    \end{itemize}
    \item Common Basis Functions: Polynomial, Fourier, Gaussian, Sigmoid
\end{itemize}

\subsubsection*{Example:}
Consider the polynomial $f(x) = \alpha_1 c + \alpha_2 x + \alpha_3 x^3 + \alpha_4 x^3$
\begin{itemize}
	\item The basis functions are $\phi_1 = c, \phi_2 = x, \phi_3 = x^3, \phi_4 = x^3$
	\item We could construct a $k \times 4$ design matrix, where the points are selected using any of the previously defined criteria (e.g., G-optimality)
	\[X =\begin{bmatrix} 1 & x_1 & x_1^2 & x_1^3 \\ 1 & x_2 & x_2^2 & x_2^3 \\ 1 & x_3 & x_3^2 & x_3^3 \\ 1 & x_4 & x_4^2 & x_4^3 \end{bmatrix}\]
    \item After obtaining labels, we can use OLS to estimate $\alpha = [\alpha_1 \alpha_2 \alpha_3 \alpha_4]'$
\end{itemize}

\subsection*{Summary: Optimal Design of Experiments}
\begin{itemize}
	\item DOE is primarily used for learning parametric regression models that can be expressed as linear basis functions
	\item DOE is similar to Active Learning in the sense that it selects training data
	\begin{itemize}
	    \item DOE methods select instances by optimizing some scalar property of $X'X$ which will, in turn, minimized the (expected) generalization error
	    \item Unlike active learning, DOE methods select the design points at the same time
    \end{itemize}
\end{itemize}

\section*{Recursive Dyadic Partition (RDP) Algorithm}
\begin{itemize}
	\item RDP is an algorithm for learning \textbf{non-parametric regression models}
	\item Type II (cluster expoiltation).  However, unlike ZLG, DH, and PLAL, it is intended for the \textbf{membership query synthesis} data access model.
	\begin{itemize}
	    \item Pool-based sampling is possible, if certain conditions are met
    \end{itemize}
    \item It is \textbf{aggressive}
    \item Guarantees:
    \begin{itemize}
	    \item RDP will converge to the optimal model
	    \item The upper and lower bounds on the rate of convergence are known
    \end{itemize}
\end{itemize}

\subsection*{Rationale}
RDP is intended for learning \textbf{piecewise smooth functions}.
\begin{center} 
	\includegraphics*[width=0.8\textwidth]{W9_1.png} 
\end{center}
For functions whose complexity is spatially non-uniform and highly concentrated in small subsets of the domain, \textbf{the extra spatial adaptivity of the active learning paradigm can lead into significant performance gains}
\begin{itemize}
	\item The authors show that \textbf{there is no benefit to active learning} (relative to passive learning), if the function satisfies the \textbf{H\"{o}lder condition}.
\end{itemize}

\subsection*{RDP: Dyadic Space-Partitioning Trees}
The algorithm uses dyadic space-partitioning trees.
\begin{itemize}
	\item Example in a 2D domain ($2^d$-ary tree structure)
	\begin{center} 
	    \includegraphics*[width=0.6\textwidth]{W9_2.png} 
    \end{center}
    \item Consider $f \::\: [0, 1]^d \rightarrow \mathbb{R}$
    \begin{itemize}
	    \item Step 1: Divide $[0, 1]^d$ into $2^d$ equal sized hypercubes
	    \item Step 2: Repeat this process again on each hypercube
    \end{itemize}
    \item Repeating this process $\log(m)$ times gives rise to a partition of the unit cube into $m^d$ hypercubes of identical size.  (This process is represented as a $2^d$-ary tree structure above.)
\end{itemize}
Note: Dyadic means to divide the domain by $2^d$, where $d$ is the dimension of $\mathcal{X} = [0, 1]^d$\\\\
There are \textbf{two phases} for this algorithm.
\begin{enumerate}
	\item \textbf{Preview Phase:}
	\begin{itemize}
	    \item Use half of the labeling budget to label points \textbf{uniformly} within the domain, $\mathcal{X}$
	    \item Learn a bounded height regression tree modified to partition the domain via a \textbf{Recursive Dyadic Partition} (RDP)
	    \item Use the RDP model to identify boundaries in the domain
    \end{itemize}
    \item \textbf{Refinement Phase:}
    \begin{itemize}
	    \item Use the remaining half of the budget to label the boundary regions densely
	    \item Learn a regression model from the labeled points
    \end{itemize}
\end{enumerate}

\subsubsection*{Example:}
\textbf{Preview Phase:}
\begin{center} 
	\includegraphics*[width=0.7\textwidth]{W9_3.png} \\
    The deepest nodes in the tree will tend to be in boundary regions (i.e., places where the function changes quickly)
\end{center}
\textbf{Refinement Phase:}
\begin{center}
    \includegraphics*[width=0.7\textwidth]{W9_4.png} \\
    Use the remaining budget to desnsely sample the regions of the domain corresponding to the deepest nodes of the regression tree
\end{center}

\subsubsection*{Technical Issue}
\begin{itemize}
	\item If the dyadic partition is aligned with a boundary, the refinement phase is unlikely to sample in that region
	\item This can be addressed by building multiple RDP models, shifted along one or more dimensions.  For example, the bottom left corner in the above examples.
\end{itemize}

\subsection*{RDP Guarantees}
\textbf{Consistency and Rate of Convergence} (under certain assumptions on $f$)
\begin{itemize}
	\item The authors derive upper and lower bounds for the rate of convergence to the optimal model
	\begin{itemize}
	    \item Upper bound: $\mathbb{E} \left[ \Vert \hat{f} - f \Vert^2\right] \leq C \left(\frac{n}{\log n}\right)^{\frac{1}{d - 1 + \frac{1}{d}}}$
	    \item Lower bound: $\mathbb{E} \left[\Vert \hat{f} - f \Vert^2\right] \leq cn^{-\frac{1}{d - 1}}$
    \end{itemize}
    \item The corresponding bounds for passive learning are as follows:
    \begin{itemize}
	    \item Upper bound: $\mathbb{E} \left[ \Vert \hat{f} - f \Vert^2\right] \leq C \left(\frac{n}{\log n}\right)^{-\frac{1}{d}}$
	    \item Lower bound: $\mathbb{E} \left[\Vert \hat{f} - f \Vert^2\right] \leq cn^{-\frac{1}{d}}$
    \end{itemize}
    \item That is, RDP's complexity is essentially that of passive learning, but on a lower-dimensional space
\end{itemize}

\subsection*{RDP Algorithm: Summary}
\begin{itemize}
	\item The algorithm is elegant and easy to implement, but \textbf{the number of samples needed to obtain reasonable accuracies will be large in practice}, since we need to guarantee sufficient uniform sampling in a $d$-dimensional space, which is expensive (see \textit{curse of dimensionality})
	\item If desired, we can use the labeled examples to train a more sophisticated regression model
	\item The algorithm is noteworthy because it is \textbf{aggressive}, yet the authors are able to prove theorems about learning rates
    \begin{itemize}
	    \item Note, however, that these guarantees only hold if the underlying function is piecewise continuous
    \end{itemize}
\end{itemize}

\subsection*{Summary: DOE vs. Active Learning}
Objective: Generate dataset from which one can learn an accurate predictive model
\begin{center}
\begin{tabular}{|c|c|c|}
\hline
\textbf{Category} & \textbf{Design of Experiments} & \textbf{Active Learning} \\ \hline
Logistics & \begin{tabular}[c]{@{}c@{}}Limited to one batch\\ (e.g., clinical trials)\end{tabular} & \begin{tabular}[c]{@{}c@{}}Multiple batches\\ (e.g., compound screening)\end{tabular} \\ \hline
Feature knowledge & \begin{tabular}[c]{@{}c@{}}Good when you know all your\\ features are useful\end{tabular} & \begin{tabular}[c]{@{}c@{}}Good when you don't know if\\ your features are useful\end{tabular} \\ \hline
Model assumptions & \begin{tabular}[c]{@{}c@{}}Typically assumes data are\\ derived from a specific model\end{tabular} & No assumptions needed \\ \hline
\end{tabular}
\end{center}

\section*{ALICE}
\textbf{ALICE:} ``Active Learning in Approximately Linear Regression Based on Conditional Expectation of Generalization Error'', Sugiyama (JMLR, 2006)

\subsection*{Aleatoric Uncertainty vs. Epistemic Uncertainty}
\begin{itemize}
	\item \textbf{Aleatoric uncertainty:} also known as data uncertainty, is irreducible uncertainty in data that gives rise to uncertainty in predictions
	\begin{itemize}
	    \item This type of uncertainty is not a property of the model, but rather is an inherent property of the data distribution, and hence, it is irreducible
    \end{itemize}
    \item \textbf{Epistemic uncertainty:} also known as knowledge uncertainty, occurs due to inadequate knowledge
    \begin{itemize}
	    \item One can define models to answer different questions in model-based prediction
    \end{itemize}
\end{itemize}

\subsection*{The ALICE Algorithm}
\textbf{Key Features:}
\begin{itemize}
	\item ALICE is an algorithm for learning parametric regression models that take the form of \textbf{linear basis functions}
	\item Unlike most query selection methods, \textbf{it is not iterative!}
	\begin{itemize}
	    \item It constructs an A-optimal design, and then queries those points all at once
	    \item It is intended for the membership query synthesis data access model
	    \begin{itemize}
	        \item Pool-based sampling is possible, if certain conditions are met
        \end{itemize}
    \end{itemize}
    \item It is \textbf{aggressive}, but corrects for sampling bias using \textbf{importance weighting} (similar to IWAL)
    \item Guarantees
    \begin{itemize}
	    \item ALICE will converge to the optimal model
	    \item Under the assumption that the true model is approximately linear, the author proves bounds on the learning rate
    \end{itemize}    
\end{itemize}

\subsection*{ALICE: Motivation}
Traditional DOE methods \textbf{assume that we know the form of the true function}
\begin{itemize}
	\item Suppose that we are wrong.  How bad could the generalization error be?
	\item For example:
    \begin{itemize}
	    \item True function: $f(x) = 1 - x + x^2 + \delta x^3 + \epsilon$
	    \item Assume that $\epsilon \sim \mathcal{N}(0, 0.3)$
	    \item Our assumed model: $g(x) = \alpha_1 - \alpha_2 x + \alpha_3 x^2$
	    \item Compare A-optimal design vs a random selection of $k = 10$ labeled instances
    \end{itemize}
    \begin{center}
        \begin{tabular}{|c|c|c|c|}
        \hline
        Generalization Error (k = 10) & $\delta = 0$ & $\delta = 0.05$ & $\delta = 0.5$ \\ \hline
        A-optimal design: $\min \text{Tr}\left((X'X)^{-1}\right)$ & 1.1 & 2.0 & 95.5 \\ \hline
        Random Selection & 2.2 & 2.3 & 7.1 \\ \hline
        \end{tabular}
    \end{center}
    \item The magnitude of the error for $\delta = 0.5$ is due in part to the bias we introduce by misspecifying the model
\end{itemize}

\subsection*{ALICE: Background}
It can be shown that for any linear estimator that:
\[Variance = \sigma^2 \text{Tr}(ULL') \propto \text{Tr}(ULL')\]
where $L = (X'X)^{-1} X'$ is the \textbf{Moore-Penrose pesudoinverse} and
\[U_{i, j} = \int \phi_i(x) \phi_i(x) P_{\text{test}}(x) \:\dd x\]
is a $b \times b$ matrix, where $b$ is the number of terms in the linear basis function
\begin{itemize}
	\item Recall that \textbf{for a correctly specified model}, ordinary least squares (OLS) is an unbiased estimator, thus, \textbf{minimizing $\text{Tr}(ULL')$ will minimize the (expected) generalization error}.
\end{itemize}
The author decomposes bias into two orthogonal components:
\begin{itemize}
	\item The \textbf{out-of-model bias} due to misspecifying the model;
	\[B_{out} = \int \left(f(x) - g(x) \right)^2 P_{\text{test}}(x) \:\dd x\]
    \item The \textbf{in-model bias} due to differences in $P_{\text{train}}$ and $P_{\text{test}}$
    \[B_{in} = \int \left(\mathbb{E}_{\text{train}} [\hat{f}(x)] - g(x)\right)^2 P_{\text{test}(x) \:\dd x}\]
\end{itemize}
\begin{center} 
	\includegraphics*[width=0.5\textwidth]{W10_1.png} 
\end{center}
\begin{itemize}
	\item If we have misspecified the model, $B_{out} > 0$
	\begin{itemize}
	    \item But, we can't reduce it either (unless we are \textbf{told} the true form of the function)
	    \begin{itemize}
	        \item That is, $B_{out}$ is a kind of irreducible error
        \end{itemize}
        \item So, we can just ignore $B_{out}$ when it comes to selecting design points
    \end{itemize}
    \item OLS is an unbiased estimator, so is it reasonable to assume that $B_{in} = 0$?
    \begin{itemize}
	    \item Unfortunately, OLS is only unbiased if $P_{train} = P_{test}$, which isn't true in Active Learning, so $B_{in} > 0$
	    \item Solution: Use \textbf{Importance Weighting} (just like the IWAL algorithm)
    \end{itemize}
\end{itemize}
ALICE's overall strategy is to use importance weighting to reduce $B_{in}$, and then to use an A-optimal design criterion to reduce variance, and thus the generalization error

\subsection*{Recall: Importance Weighting}
\begin{itemize}
	\item The empirical average loss for model $h$ on $n$ instances is: $\frac{1}{n} \sum_i loss(h(x_i), y_i)$
	\item As $n$ goes to $\infty$ the training loss converges to
	\[\int loss(h(x), y) P_{train}(x) \:\dd x\]
    \item Of course, what we really care about is the \textbf{test} loss:
    \[\int loss(h(x), y) P_{test}(x) \:\dd x\]
    \begin{itemize}
	    \item If $P_{train} = P_{test}$, then these losses will be the same
	    \item But if we have a biased training data, then $P_{train}(\mathcal{X}) \neq P_{test}(\mathcal{X})$
    \end{itemize}
    \item Recall that many learning algorithms can use \textbf{weighted loss functions}
    \begin{itemize}
	    \item Ex. weighted empirical loss: $\widehat{err}(h) = \frac{1}{n} \sum_i w_i loss(h(x_i), y_i)$
    \end{itemize}
    \item Let the \textbf{importance weight} of a training instance be defined as:
    \[w_i = \frac{P_{test}(x_i)}{P_{train}(x_i)}\]
    \item The importance-weighted loss is
    \[\widehat{err}_{IW}(h) = \frac{1}{n} \sum_i \frac{P_{test}(x_i)}{P_{train}(x_i)} loss(h(x_i), y_i)\]
    \item As $n$ goes to $\infty$ the importance-weighted loss converges to:
    \begin{align*}
        \widehat{err}_{IW}(h) &= \int loss(h(x), y) \frac{P_{test}(x)}{P_{train}(x)} P_{train}(x) \:\dd x \\
        &= \int loss(h(x), y) P_{test}(x) \:\dd x
    \end{align*}
    which is exactly what we want!
    \item Thus, if we can compute suitable importance weights, the fact that our training data are biased will not matter (in the limit)
\end{itemize}

\subsection*{ALICE and Importance Weighting}
\begin{itemize}
	\item Recall the \textbf{oridinary least squares} (OLS) estimate for a linear model $y = X\beta$ is:
	\[\hat{\beta} = (X'X)^{-1} X' y\]
    \item The \textbf{importance-weighted Least Squares} estimate minimizes the weighted sum of errors
    \[\hat{\beta} = (X' DX)^{-1} X' Dy\]
    where 
    \[D_{i, j} = \begin{cases} \frac{P_{test}(x_i)}{P_{train}(x_i)} & \text{if } i = j \\ 0 & \text{otherwise} \end{cases}\]
\end{itemize}
Properties of \textbf{Importance-weighted Least Squares}
\begin{itemize}
	\item It is unbiased, asymptotically, even if the model is misspecified!
	\item If the model is approximately correct (i.e., $f(x) \approx g(x)$), then the in-model bias is much smaller than the variance (i.e., $B_{in} << \text{Variance}$)
	\item If we re-define matrix $L \equiv (X' D X)^{-1} X' D$, then the variance is still
	\[Variance = \sigma^2 \text{Tr}(ULL') \propto \text{Tr}(ULL')\]
    \begin{itemize}
	    \item So, we can still use an A-optimal design criterion to reduce variance, and thus the generalization error
    \end{itemize}
\end{itemize}

\subsection*{ALICE: Implementation Details}
\begin{itemize}
	\item The traditional approach to finding A-optimal (or D-optimal, etc.) designs is to use \textbf{semi-definite programming}
	\item ALICE doesn't use semi-definite programming
	\begin{itemize}
	    \item Instead, it takes as input a finite set of $k$ probability densities, $\hat{P} = \{P_1, \cdots, P_k\}$
	    \item It \textbf{samples} $n$ instances from each $p \in \hat{P}$ and then
	    \begin{itemize}
	        \item Constructs $k$ design matrices, $X_1, \cdots, X_k$, one for each $p \in \hat{P}$
	        \item Constructs the corresponding $L_1, \cdots, L_k$ matrices
        \end{itemize}
        \item The algorithm finds the matrix $L_i$ that minimizes $\text{Tr}(UL_i L_i')$ and then requests for the corresponding set of points
        \item It then learns the regression model using importance-weighted least squares
    \end{itemize}
\end{itemize}

\subsection*{The ALICE Algorithm}
\begin{itemize}
	\item \textbf{Input:} A finite set $\hat{\mathcal{P}}$ of strictly positve probability densities
	\item Calculate $U$.  
	\[U_{i, j} = \int \phi_i(x) \phi_j(x) P_{test} (x) \:\dd x\]
    \item For each $p \in \hat{\mathcal{P}}$
    \begin{itemize}
	    \item Create training input points $\{x_i^{(p)}\}_{i = 1}^n$ following $p(x)$
	    \item Calculate $L_W$
	    \[L = (X' DX)^{-1} X' D\]
        \item Calculate $J(p)$
        \[J(p) = \text{Tr}(ULL')\]
    \end{itemize}
    \item Choose $\hat{p}$ that minimizes $J$
    \item Put $x_i = x_i^{(\hat{p})}$ for $i = 1, 2, \dots, n$.
    \item Call the oracle and learn the model $\alpha = Ly$, which entails:
    \begin{itemize}
        \item Observe the training output values $\{y_i\}_{i = 1}^n$ at $\{x_i\}_{i = 1}^n$
        \item Calculate $\hat{\alpha}_W$ by Eq. (9) 
    \end{itemize}
    \item \textbf{Output:} $\hat{\alpha}_W$
\end{itemize}

\subsection*{ALICE Guarantees}
\textbf{Error} (assuming that $f(x) = g(x) + \delta r(x)$)
\begin{itemize}
	\item $\mathbb{E}[Err(x)] = \sigma^2 \text{Tr}(ULL') + O(\delta n^{-1})$
    \begin{itemize}
	    \item Note that $\sigma^2 \text{Tr}(ULL') = O(n^{-1})$, so if $\delta = O(1)$, then the expeted error decreases as $O(n^{-1})$ which is the same asymptotic rate as if the model was specified correctly
    \end{itemize}
\end{itemize}

\subsection*{Example Results}
\begin{itemize}
	\item True function: $f(x) = 1 - x + x^2 + \delta x^3 + \epsilon$
	\begin{itemize}
	    \item Assume that $\epsilon \sim \mathcal{N}(0, 0.3)$
    \end{itemize}
    \item Our assumed model: $g(x) = \alpha_1 - \alpha_2 x + \alpha_3 x^2$
    \item Compare three different strategies using $k = 10$ labeled instances
\end{itemize}
\begin{center}
    \begin{tabular}{|c|c|c|c|}
    \hline
    Generalization Error (k = 10) & $\delta = 0$ & $\delta = 0.05$ & $\delta = 0.5$ \\ \hline
    ALICE & 1.1 & 1.2 & 4.7 \\ \hline
    A-optimal design: $\min \text{Tr}\left((X'X)^{-1}\right)$ & 1.1 & 2.0 & 95.5 \\ \hline
    Random Selection & 2.2 & 2.3 & 7.1 \\ \hline
    \end{tabular}
\end{center}

\subsection*{ALICE: Summary}
\begin{itemize}
	\item ALICE actively learns \textbf{linear basis function models}, and provides generalization guarantees, \textbf{even if the model is misspecified}
	\item It is intended for a membership query synthesis data access model
	\item It uses two techniques to (approximately) minimize generalization error
	\begin{itemize}
	    \item \textbf{Importance weighting} corrects for in-model bias, resulting in an unbiased estimator
	    \item The importance weighted \textbf{A-optimal design} reduces variance
    \end{itemize}
    \item It has the same learning rate as passive learning, but produces better models
    \item \textbf{It assumes that we know $P_{test}$}
    \begin{itemize}
	    \item If you don't know $P_{test}$ but do have access to an unlabeled pool of samples from $P_{test}$, the same author does have a closely related algorithm: 
        \begin{itemize}
	        \item \underline{Pool-based active learning in approximate linear regression}
        \end{itemize}
    \end{itemize}
\end{itemize}

\section*{LapRDD}
``Laplacian Regularized D-optimal Design for Active Learning and its Application to Image Retrieval'', He (IEEE Trans Image Process, 2009)




\end{document}



